% ============================================================
% Section 7/8: Conclusion & Future Work
% Target Journal: Neural Computing and Applications (NCA)
% Master Guides: NCA_WRITING_DNA.md, SCIENTIFIC_REASONING_DNA.md
% ============================================================

\section{Conclusion \& Future Work}\label{sec:conclusion}

This paper presented a controlled benchmark study evaluating five metaheuristic and stochastic optimization algorithms—Random Search (RS), Genetic Algorithm (GA), Particle Swarm Optimization (PSO), Differential Evolution (DE), and Grey Wolf Optimizer (GWO)—for hyperparameter selection in neural intraday trend classification. By enforcing an identical computational budget of $N_{\text{eval}} = 1,500$ fitness evaluations per seed, a leakage-free 60/20/20 chronological temporal split, and a compound MCC-driven fitness objective ($0.60\text{MCC} + 0.40F_1$), the benchmark eliminated common evaluation biases prevalent in published financial machine learning literature.

The empirical findings yield three primary conclusions:
\begin{enumerate}
    \item \textbf{Guided Metaheuristics Significantly Outperform Unguided Search}: All four metaheuristics achieved statistically significant predictive improvements over the Random Search baseline ($p < 0.05$, Cohen's $d > 0.80$). GA achieved the highest out-of-sample MCC ($0.231 \pm 0.012$), representing a relative performance gain of +24.8\% over Random Search ($0.185 \pm 0.021$).
    \item \textbf{Recombination Reconciles Exploration and Exploitation}: Evolutionary paradigms with strong crossover mechanisms (GA and DE) exhibited higher long-term robustness than leader-driven swarms (GWO), avoiding premature local convergence in multimodal hyperparameter landscapes.
    \item \textbf{Predictive Gains Translate into Real Trading Profitability}: In out-of-sample financial backtesting with realistic transaction fees ($0.01\%$ per trade), the GA-optimized model generated a Net Return of $+18.4\%$ and an annualized Sharpe Ratio of $1.42$, confirming that compound MCC optimization delivers genuine economic utility beyond raw accuracy.
\end{enumerate}

Future research will expand this benchmark framework along three key avenues:
\begin{enumerate}
    \item \textbf{Multi-Asset Universe Generalization}: Extending the temporal benchmark from B3 Mini-Index futures (WIN) to a broader multi-asset universe including Full Index (IND), Mini-Dollar (WDO), and Full Dollar (DOL) contracts over sub-minute tick resolutions.
    \item \textbf{Hybrid Swarm-Evolutionary Operators}: Investigating hybrid metaheuristics (e.g., GA-PSO hybrids) that combine PSO's rapid early exploitation with GA's late-stage diversification mechanics.
    \item \textbf{Pareto Multi-Objective Optimization}: Formulating hyperparameter selection as a Pareto multi-objective problem to simultaneously optimize predictive MCC, Sharpe Ratio, and computational execution latency.
\end{enumerate}
